\section{Resultados por incremento}
El desarrollo produjo un framework funcional cuyos cuatro incrementos están
consolidados y respaldados por pruebas automatizadas.

\subsection{Incremento 1: API backend y modelado (consolidado)}
Se dispone de un entorno de despliegue reproducible con \texttt{docker
compose}, probado en sistemas operativos y arquitecturas de CPU distintas, y
de un núcleo de modelos ORM organizado por dominio, con herencia polimórfica
para autenticación/autorización, prefijo de tablas configurable, versionado
con \gls{continuum} y migraciones con Alembic. El punto de entrada
\texttt{initialize(config)} permite usar el backend como librería.

\subsection{Incremento 2: contrato back--front (consolidado)}
El contrato uniforme está implementado y documentado: envoltorio de respuesta
normalizado, manejadores de error globales, catálogo de códigos estándar,
fábricas CRUD, endpoints de descubrimiento y una capa cliente genérica tipada
en Angular que convive con el código heredado sin romperlo. Las pruebas
\texttt{test\_envelope.py} y \texttt{test\_auth.py} verifican la forma del
envoltorio, los errores y el descubrimiento.

\subsection{Incremento 3: formularios dinámicos (consolidado)}
La cadena esquema a formulario funciona extremo a extremo: el
backend enriquece el esquema con pistas \texttt{x-*} a partir de los tipos
SQLAlchemy y el frontend genera, sin HTML específico, un formulario Formly
completo con validadores y expresiones condicionales. El componente
reutilizable \texttt{<ngt-dynamic-form>} encapsula el flujo completo
(descubrimiento--esquema--formulario--envío) y las pruebas 
\texttt{test\_form\_schema.py} y \texttt{schema-to-formly.spec.ts} cubren los casos relevantes.

\subsection{Incremento 4: bundle de esquemas}
El \emph{registry} de entidades, el endpoint \texttt{/api/sys/schemas}, el
versionado por ETag y sus pruebas (\texttt{test\_schema\_bundle.py}) están
implementados y operativos. El incremento se completó añadiendo las dos piezas
que quedaban pendientes: el reflejo del \textbf{polimorfismo} de
\texttt{FunctionalObject} y subclases mediante \texttt{oneOf} con un
\texttt{x-discriminator} sobre \texttt{object\_type\_id}, y el de los
\texttt{relationship()} del ORM como extensión \texttt{x-relationships}. Ambas
viajan dentro del propio esquema enriquecido, son deterministas (no
desestabilizan el ETag del \emph{bundle}) y conservan las \texttt{properties}
base, de modo que un consumidor que no entienda \texttt{oneOf} sigue
funcionando. Las pruebas verifican que la entidad raíz expone
\texttt{oneOf}+\texttt{x-discriminator} con sus subclases, que una subclase no
lo emite, y que \texttt{x-relationships} refleja las relaciones resueltas
contra el \emph{registry}.

\section{Valoración global}
La tabla~\ref{tab:result} resume el estado de los objetivos específicos. El
resultado más significativo es cualitativo: exponer y editar una entidad nueva
pasa de requerir un router, esquemas y formularios escritos a mano a
\emph{declarar el modelo ORM} y dejar que el framework derive el resto. Esto
resuelve el objetivo de partida sobre un modelo de datos en evolución.

\subsection{Comparativa antes/después}
La forma más clara de valorar el impacto es comparar el trabajo necesario para
poner en producción una entidad nueva, editable desde el frontend, antes y
después del framework (tabla~\ref{tab:antesdespues}).

\begin{table}[!htbp]
\caption{Trabajo para exponer y editar una entidad nueva.}
\label{tab:antesdespues}
\centering
\begin{tabular}{|p{4.4cm}|p{4.4cm}|p{4.4cm}|}
\hline
\textbf{Tarea} & \textbf{Antes} & \textbf{Después} \\ \hline
Endpoints CRUD & Router escrito a mano & Una fábrica CRUD \\ \hline
Validación & Esquemas Pydantic a mano & Derivada del ORM \\ \hline
Contrato/errores & \emph{Ad hoc} por endpoint & Envoltorio uniforme \\ \hline
Cliente frontend & Método dedicado en \texttt{BackendService} &
\texttt{EntityClient.for(path)} \\ \hline
Formulario & HTML + lógica por campo & \texttt{<ngt-dynamic-form
entityPath>} \\ \hline
Coste ante cambio del modelo & Propagación manual a 3 capas & Automática
desde el ORM \\ \hline
\end{tabular}
\end{table}

El cambio cualitativo es el relevante: el esfuerzo deja de crecer con el
número de entidades y con la frecuencia de cambios del modelo, que era
precisamente el problema de partida. La verificación de que la cadena no se
rompe ante un cambio del modelo queda garantizada por las pruebas de la
heurística de pistas y del conversor.

\begin{table}[!htbp]
\caption{Grado de consecución de los objetivos específicos.}
\label{tab:result}
\centering
\begin{tabular}{|p{8.5cm}|c|}
\hline
\textbf{Objetivo específico} & \textbf{Estado} \\ \hline
Crear una API modular & Completado \\ \hline
Diseñar y modelar los esquemas de datos & Completado \\ \hline
Comunicación estandarizada back--front & Completado \\ \hline
Generador automático de formularios (Formly) & Completado \\ \hline
Generación automática de esquemas desde el ORM & Completado \\ \hline
\end{tabular}
\end{table}

\subsection{Métricas cuantitativas}
Más allá de la valoración cualitativa, el resultado admite una caracterización
cuantitativa. La tabla~\ref{tab:metricas} recoge las cifras más relevantes,
obtenidas directamente del código fuente y de la ejecución de la batería de
pruebas.

\begin{table}[!htbp]
\caption{Métricas cuantitativas del framework resultante.}
\label{tab:metricas}
\centering
\begin{tabular}{|p{9.5cm}|c|}
\hline
\textbf{Métrica} & \textbf{Valor} \\ \hline
Entidades expuestas por la cadena ORM--esquema & 61 \\ \hline
Dominios funcionales (más \emph{plugins}) & 9 \\ \hline
Endpoints CRUD generados automáticamente por entidad & 6 \\ \hline
Endpoints CRUD generados sin escribir a mano ($61\times6$) & 366 \\ \hline
Tipos de \gls{jsonschema} mapeados a Formly & 16 \\ \hline
Extensiones de esquema \texttt{x-*} reconocidas & 39 \\ \hline
Tests automatizados de backend (\emph{pytest}) & 41 \\ \hline
Cobertura de los tests de backend & 45\,\% \\ \hline
Tests unitarios del núcleo esquema--Formly & 17 \\ \hline
\end{tabular}
\end{table}

Las 61 entidades, repartidas en nueve dominios funcionales, se exponen sin
escribir manualmente sus operaciones: la fábrica CRUD genera seis \emph{endpoints}
por entidad, frente a los más de doscientos métodos escritos a mano que
requería el backend monolítico original. La cadena de generación reconoce
dieciséis tipos de datos y treinta y nueve extensiones \texttt{x-*} de
presentación, permisos y relaciones. La batería de pruebas automatizadas del
backend (41 pruebas, con un 45\,\% de cobertura) y los tests unitarios del
conversor esquema--Formly respaldan la estabilidad de la cadena ante cambios
del modelo.
